
%%%%%%%%%%%%%%%%%%%%%%%%%%%%%%%%%%%%%%%%%%%%%%%%%%%%%%%%%%%
%% Inicio del capítulo de análisis del problema
%%%%%%%%%%%%%%%%%%%%%%%%%%%%%%%%%%%%%%%%%%%%%%%%%%%%%%%%%%%

\chapter{Análisis del problema}
\label{cap:analisisProblema}

%% El paquete babel (spanish) activa "<" y ">" como atajos para «
%% y », lo que interfiere con la sintaxis de flechas de TikZ
%% (p. ej. ->, <->). Se desactivan dichos atajos durante el capítulo.
\shorthandoff{<>}

En este capítulo se aborda el  problema  definido en el presente  Trabajo Fin de Máster, centrado en la integración de LLM en el modelado de procesos de negocio mediante \textit{bpmn.io} y en la incorporación de un módulo de análisis del valor percibido por el cliente basado en la clasificación PERVAL.

 Con el propósito de comprender el problema y fundamentar el diseño de la solución propuesta, se realiza, en primer lugar,   un modelado funcional del sistema mediante diagramas de actores, de contexto y de casos de uso. Estos modelos permiten delimitar el alcance del sistema, identificar sus principales componentes y especificar  las interacciones previstas entre el usuario, la extensión sobre \textit{bpmn.io} y los proveedores de LLM. 

Posteriormente, se  analiza el marco legal y ético asociado al uso de modelos de lenguaje y al tratamiento de los datos introducidos por  los usuarios. Asimismo,  se  identifican y comparan las principales alternativas tecnológicas consideradas para abordar el problema  planteado, justificando las decisiones  de diseño finalmente adoptadas. 


 Finalmente, a partir de dicho análisis realizado, se concreta la solución propuesta mediante la especificación de los requisitos funcionales y no funcionales del sistema, la definición de  un conjunto de historias de usuario y la planificación del trabajo realizado, organizada en \textit{sprints} semanales siguiendo una metodología ágil  basada en Scrum.

%%%%%%%%%%%%%%%%%%%%%%%%%%%%%%%%%%%%%%%%%%%%%%%%%%%%%%%%%%%
\section{Modelado funcional del sistema}
\label{sec:modeladoFuncional}

Con el objetivo de delimitar de forma precisa el alcance funcional de l{a solución propuesta y comprender las interacciones entre sus diferentes componentes, se ha realizado un modelado funcional del sistema utilizando basado en diagramas UML simplificados. El modelado funcional constituye una actividad fundamental durante el análisis de sistemas, ya que permite identificar los actores involucrados, especificar las responsabilidades de cada componente y describir las funcionalidades que el sistema debe proporcionar para satisfacer las necesidades de los usuarios.

En el contexto del presente TFM, el modelado funcional resulta especialmente relevante debido a la naturaleza híbrida de la solución desarrollada, que integra una extensión sobre bpmn.io, un servidor intermedio basado en Model Context Protocol (MCP) y uno o varios proveedores externos de modelos de LLMs. La interacción coordinada entre estos componentes requiere definir de manera explícita las fronteras del sistema, las responsabilidades de cada subsistema y los flujos de información intercambiados entre ellos.

Para ello, el análisis se estructura en tres niveles complementarios. En primer lugar, se identifican los actores que participan en el sistema y las relaciones existentes entre ellos mediante un diagrama de actores. En segundo lugar, se presentan los diagramas de contexto, cuyo objetivo es delimitar el alcance de cada subsistema y representar las interacciones establecidas con los actores identificados. Finalmente, se describen los diagramas de casos de uso, que permiten especificar las funcionalidades ofrecidas por cada subsistema y la forma en que dichas funcionalidades son utilizadas para satisfacer los objetivos del usuario.

%%%%%%%%%%%%%%%%%%%%%%%%%%%%%%%%%%%%%%%%%%%%%%%%%%%%%%%%%%%
\subsection{Diagrama de actores}
\label{subsec:diagramaActores}

En el sistema propuesto se identifican dos actores principales, representados en la figura~\ref{fig:diagramaActores}:

\begin{itemize}
    \item \textbf{Usuario}: persona que utiliza la extensión sobre \textit{bpmn.io} para describir, generar, editar y analizar modelos de procesos de negocio en notación BPMN. Puede ser un analista de procesos, un consultor o cualquier perfil que necesite modelar procesos sin un conocimiento profundo de la notación BPMN.
    \item \textbf{LLM}: servicio externo de modelos de lenguaje (\textit{GitHub Models}) utilizado por el sistema para interpretar las descripciones en lenguaje natural del usuario y generar o modificar los modelos BPMN, así como para realizar el análisis de valor PERVAL.
\end{itemize}

Los componentes internos de la solución son la extensión sobre bpmn.io y el servidor MCP y se representan dentro de la frontera del sistema.

\begin{figure}[h!]
\centering
\begin{tikzpicture}[scale=1.0, >=Stealth, line width=0.8pt]

% Frontera del sistema (componentes internos)
\draw[thick, dashed] (2.0,-0.8) rectangle (7.0,1.8);
\node[font=\bfseries] at (4.5,2.1) {Sistema};
\node[align=center, font=\small] at (4.5,1.1) {Extensión \textit{bpmn.io}};
\node[align=center, font=\small] at (4.5,0.2) {Servidor MCP};

% Actor externo: Usuario (izquierda)
\begin{scope}[shift={(-1.5,0)}]
  \stickman
  \node[below=0.15cm] at (0,-0.8) {Usuario};
\end{scope}

% Actor externo: Proveedor de LLM (derecha)
\begin{scope}[shift={(10.5,0)}]
  \stickman
  \node[below=0.15cm, align=center] at (0,-0.8) {Proveedor\\de LLM};
\end{scope}

% Relaciones actor-sistema
\draw[<->] (-0.3,0) -- (2.0,0);
\draw[<->] (7.0,0) -- (9.8,0);

\end{tikzpicture}
\caption{Diagrama de actores del sistema.}
\label{fig:diagramaActores}
\end{figure}

%%%%%%%%%%%%%%%%%%%%%%%%%%%%%%%%%%%%%%%%%%%%%%%%%%%%%%%%%%%
\label{subsec:diagramasContexto}

Para delimitar con claridad el alcance de cada subsistema y sus relaciones con los actores identificados en la sección anterior, se presentan a continuación dos diagramas de contexto: el de la extensión cliente sobre \textit{bpmn.io} y el del servidor MCP que actúa de \textit{backend}.


La figura~\ref{fig:contextoExtension} muestra el contexto de la \textbf{extensión sobre \textit{bpmn.io}}, que constituye la interfaz con la que interactúa el \textbf{Usuario}. Esta extensión se comunica con el \textbf{Servidor MCP} (componente interno del sistema) para delegar en él la generación y edición de modelos BPMN, así como el análisis PERVAL.


\begin{figure}[h!]
\centering
\begin{tikzpicture}[
  sysbox/.style={draw, thick, rectangle, minimum width=5.2cm, minimum height=2cm, align=center, fill=gray!8},
  extbox/.style={draw, thick, rectangle, minimum width=3.6cm, minimum height=1.6cm, align=center},
  >=Stealth, line width=0.8pt
]

\node[sysbox] (cliente) at (0,0) {Extensión \textit{bpmn.io}\\(cliente)};

\begin{scope}[shift={(-5,0)}]
  \stickman
  \node[below=0.2cm] at (0,-0.8) {Usuario};
\end{scope}

\node[extbox] (mcp) at (5,0) {Servidor MCP\\(\textit{backend})};

\draw[<->] (-3.6,0.3) -- (-2.6,0.3);
\draw[<->] (2.6,0) -- (3.2,0);

\end{tikzpicture}
\caption{Diagrama de contexto de la extensión sobre \textit{bpmn.io} (cliente).}
\label{fig:contextoExtension}
\end{figure}

La figura~\ref{fig:contextoMCP} muestra el contexto del \textbf{servidor MCP}, que actúa como puente entre la extensión cliente y los proveedores de \textbf{LLM}. Este componente recibe las peticiones de generación o edición de modelos BPMN procedentes de la extensión, construye los \textit{prompts} adecuados, invoca al LLM seleccionado y devuelve el XML BPMN resultante, validado y, en caso necesario, reparado.

\begin{figure}[h!]
\centering
\begin{tikzpicture}[
  sysbox/.style={draw, thick, rectangle, minimum width=5.2cm, minimum height=2cm, align=center, fill=gray!8},
  extbox/.style={draw, thick, rectangle, minimum width=3.6cm, minimum height=1.6cm, align=center},
  >=Stealth, line width=0.8pt
]

\node[sysbox] (mcp) at (0,0) {Servidor MCP\\(\textit{backend})};

\node[extbox] (cliente) at (-5,0) {Extensión \textit{bpmn.io}\\(\textit{Sistema}, cliente)};

\node[extbox] (llm) at (5,0) {LLM\\(\textit{GitHub Models}};

\draw[<->] (-3.2,0) -- (-2.6,0);
\draw[<->] (2.6,0) -- (3.2,0);

\end{tikzpicture}
\caption{Diagrama de contexto del servidor MCP (\textit{backend}).}
\label{fig:contextoMCP}
\end{figure}

%%%%%%%%%%%%%%%%%%%%%%%%%%%%%%%%%%%%%%%%%%%%%%%%%%%%%%%%%%%
\subsection{Diagramas de casos de uso}
\label{subsec:diagramasCasosUso}

Una vez delimitado el contexto de cada subsistema, se detallan a continuación los casos de uso correspondientes. La figura~\ref{fig:casosUsoExtension} muestra los casos de uso de la \textbf{extensión sobre \textit{bpmn.io}}, en la que el actor \textbf{Usuario} interactúa con las funcionalidades de descripción y edición conversacional del modelo, visualización y exportación del diagrama, configuración del análisis PERVAL y selección de idioma. Tanto la descripción inicial del proceso como la solicitud de modificaciones incluyen (\textit{\textless\textless include\textgreater\textgreater}) el caso de uso interno ``Renderizar diagrama y actualizar historial'', que el propio sistema ejecuta de forma automática. Si dicho proceso falla, se extiende (\textit{\textless\textless extend\textgreater\textgreater}) con el caso de uso ``Mostrar estado de carga o error''.

\begin{figure}[h!]
\centering
\resizebox{\linewidth}{!}{%
\begin{tikzpicture}[
  uc/.style={draw, ellipse, minimum width=3.6cm, minimum height=1.1cm, align=center, font=\small},
  >=Stealth, line width=0.7pt
]

\draw[thick] (-0.5,-8.6) rectangle (16.5,0.8);
\node[font=\bfseries] at (8.0,0.5) {Extensión \textit{bpmn.io}};

\node[uc] (uc1) at (2.4,0)    {Describir proceso en\\lenguaje natural};
\node[uc] (uc2) at (2.4,-1.5) {Solicitar modificación\\del modelo};
\node[uc] (uc3) at (2.4,-3)   {Visualizar diagrama\\BPMN};
\node[uc] (uc4) at (2.4,-4.5) {Exportar modelo\\BPMN};
\node[uc] (uc5) at (2.4,-6)   {Configurar y visualizar\\análisis PERVAL};
\node[uc] (uc6) at (2.4,-7.5) {Cambiar idioma de\\la interfaz};

\node[uc] (uc7) at (13.5,-0.75) {Renderizar diagrama y\\actualizar historial};
\node[uc] (uc8) at (13.5,-3.75) {Mostrar estado de\\carga o error};

\begin{scope}[shift={(-1.6,-3.75)}]
  \stickman
  \node[below=0.2cm] at (0,-0.8) {Usuario};
\end{scope}

\foreach \i in {1,...,6}{
  \draw (-0.9,-3.75) -- (uc\i.west);
}

\draw[dashed,->] (uc1.east) to[bend left=20] node[pos=0.5, font=\scriptsize, fill=white, inner sep=1pt] {\textless\textless include\textgreater\textgreater} (uc7.west);
\draw[dashed,->] (uc2.east) to[bend right=20] node[pos=0.5, font=\scriptsize, fill=white, inner sep=1pt] {\textless\textless include\textgreater\textgreater} (uc7.west);
\draw[dashed,->] (uc8.north) to[bend right=25] node[midway, right, font=\scriptsize, fill=white, inner sep=1pt] {\textless\textless extend\textgreater\textgreater} (uc2.east);

\end{tikzpicture}}
\caption{Diagrama de casos de uso de la extensión sobre \textit{bpmn.io}.}
\label{fig:casosUsoExtension}
\end{figure}

La figura~\ref{fig:casosUsoMCP} muestra los casos de uso del \textbf{servidor MCP}. La \textbf{extensión bpmn.io} (en su rol de cliente MCP) invoca dos casos de uso principales: ``Generar o editar modelo BPMN'' y ``Analizar valor PERVAL''. El primero incluye la construcción del \textit{prompt} de generación y la invocación al LLM, y se extiende con la limpieza y reparación del XML cuando la respuesta del modelo es incompleta o malformada (función \texttt{cleanXML}). De forma análoga, ``Analizar valor PERVAL'' incluye la construcción de un \textit{prompt} específico y la invocación al LLM para obtener la clasificación de cada tarea o entregable según las dimensiones de valor. En ambos casos, la invocación al LLM puede extenderse con la selección de un modelo de respaldo (\textit{fallback}) cuando el modelo principal no está disponible o devuelve un error.


\begin{figure}[h!]
\centering
\resizebox{\linewidth}{!}{%
\begin{tikzpicture}[
  uc/.style={draw, ellipse, minimum width=3.4cm, minimum height=1.1cm, align=center, font=\small},
  >=Stealth, line width=0.7pt
]

\draw[thick] (-0.5,-5.5) rectangle (16.5,3.6);
\node[font=\bfseries] at (8.0,3.3) {Servidor MCP};

\node[uc] (uc5) at (1.9,2.7)  {Limpiar y reparar\\XML};
\node[uc] (uc1) at (1.9,0.5)  {Generar o editar\\modelo BPMN};
\node[uc] (uc6) at (1.9,-3.7) {Analizar valor\\PERVAL};

\node[uc] (uc2) at (7.0,1.7)  {Construir prompt de\\generación};
\node[uc] (uc7) at (7.0,-2.5) {Construir prompt\\PERVAL};

\node[uc] (uc3) at (13.0,0.5)  {Invocar LLM\\(generación)};
\node[uc] (uc4) at (13.0,-1.7) {Seleccionar modelo\\de respaldo};
\node[uc] (uc8) at (13.0,-3.7) {Invocar LLM\\(análisis PERVAL)};

\begin{scope}[shift={(-2.2,-1.5)}]
  \stickman
  \node[below=0.2cm, align=center, font=\small] at (0,-0.9) {Extensión\\\textit{bpmn.io}};
\end{scope}

\begin{scope}[shift={(18.0,-1.5)}]
  \stickman
  \node[below=0.2cm] at (0,-0.8) {LLM};
\end{scope}

\draw (-1.5,-1.5) -- (uc1.west);
\draw (-1.5,-1.5) -- (uc6.west);
\draw (uc3.east) -- (15.8,-1.5);
\draw (uc8.east) -- (15.8,-1.5);

\draw[dashed,->] (uc1.east) -- node[pos=0.35, above, sloped, font=\scriptsize, fill=white, inner sep=1pt] {\textless\textless include\textgreater\textgreater} (uc2.west);
\draw[dashed,->] (uc1.east) -- node[pos=0.65, below, sloped, font=\scriptsize, fill=white, inner sep=1pt] {\textless\textless include\textgreater\textgreater} (uc3.west);
\draw[dashed,->] (uc6.east) -- node[pos=0.35, above, sloped, font=\scriptsize, fill=white, inner sep=1pt] {\textless\textless include\textgreater\textgreater} (uc7.west);
\draw[dashed,->] (uc6.east) -- node[pos=0.65, below, sloped, font=\scriptsize, fill=white, inner sep=1pt] {\textless\textless include\textgreater\textgreater} (uc8.west);

\draw[dashed,->] (uc5.south) -- node[midway, right, font=\scriptsize, fill=white, inner sep=1pt] {\textless\textless extend\textgreater\textgreater} (uc1.north);
\draw[dashed,->] (uc4.north) -- node[midway, right, font=\scriptsize, fill=white, inner sep=1pt] {\textless\textless extend\textgreater\textgreater} (uc3.south);

\end{tikzpicture}}
\caption{Diagrama de casos de uso del servidor MCP}
\label{fig:casosUsoMCP}
\end{figure}

%%%%%%%%%%%%%%%%%%%%%%%%%%%%%%%%%%%%%%%%%%%%%%%%%%%%%%%%%%%
\section{Análisis del marco legal y ético}
\label{sec:marcoLegalEtico}

El uso de LLMs para generar y modificar modelos de procesos de negocio, así como para analizar el valor percibido por el cliente a partir de información introducida por el usuario, plantea una serie de consideraciones legales y éticas que deben analizarse antes de definir la solución técnica. En esta sección se analizan tres aspectos: la protección de datos personales, la propiedad intelectual y las licencias del software empleado, y las consideraciones éticas derivadas del uso de LLM.

%%%%%%%%%%%%%%%%%%%%%%%%%%%%%%%%%%%%%%%%%%%%%%%%%%%%%%%%%%%
\subsection{Protección de datos}
\label{subsec:proteccionDatos}

El sistema desarrollado permite al usuario describir, en lenguaje natural, procesos de negocio que pueden incluir referencias a roles, departamentos o, en casos puntuales, nombres de personas u organizaciones. Estas descripciones se envían, junto con el contexto de la conversación, al proveedor de LLM seleccionado (\textit{GitHub Models}) para su procesamiento. Por tanto, resulta necesario analizar dicho tratamiento a la luz del Reglamento (UE) 2016/679, de Protección de Datos (RGPD)~\cite{rgpd2016}.

En primer lugar, cabe señalar que el sistema no requiere registro ni autenticación de usuarios, por lo que no se almacena de forma persistente ningún dato personal identificativo (nombre, correo electrónico, etc.) en los servidores de la aplicación. El servidor MCP actúa como un componente \textit{stateless}: recibe una petición, la reenvía al proveedor de LLM correspondiente, y devuelve la respuesta sin persistirla en una base de datos propia. El historial de la conversación se mantiene únicamente en el lado del cliente (en el navegador del usuario), por lo que su gestión depende del propio usuario.

No obstante, los textos introducidos por el usuario sí se transmiten a los proveedores externos de LLM, que actúan como encargados del tratamiento conforme a sus respectivas políticas de privacidad. A efectos del RGPD~\cite{rgpd2016}, la herramienta informa al usuario de que las descripciones introducidas serán enviadas a un proveedor externo de LLM para su procesamiento, de modo que el uso del sistema presupone la aceptación de estas condiciones de funcionamiento. Esta aproximación resulta más prudente que invocar un consentimiento implícito, ya que el RGPD (artículo 4.11) exige que el consentimiento sea libre, específico, informado e inequívoco, requisitos cuya acreditación formal excede el alcance de este trabajo.


%%%%%%%%%%%%%%%%%%%%%%%%%%%%%%%%%%%%%%%%%%%%%%%%%%%%%%%%%%%
\subsection{Propiedad intelectual y licencias}
\label{subsec:propiedadIntelectual}

El sistema desarrollado se apoya en distintas librerías y servicios de terceros, cuyas licencias deben ser compatibles con el uso académico y, potencialmente, con una futura distribución del código como software de libre acceso. La extensión cliente se construye sobre \textit{bpmn.io} (\textit{bpmn-js} y \textit{diagram-js}), distribuido bajo licencia \textit{bpmn.io license} (basada en MIT con restricciones sobre la eliminación del logotipo), mientras que el servidor MCP se implementa con \textit{Express} y el \textit{SDK} oficial de \textit{Model Context Protocol}~\cite{anthropic2024mcp}, ambos de licencia \textit{MIT}. El uso de los servicios de \textit{GitHub Models} está sujeto a su respectivo términos de servicio y políticas de uso de API, que deben respetarse en cuanto a límites de uso, atribución y restricciones sobre el contenido generado.

Por otra parte, los modelos BPMN generados a partir de las descripciones del usuario son obra derivada de la entrada proporcionada por este, por lo que la propiedad intelectual de dichos modelos corresponde al usuario que los genera, sin que el sistema reclame ningún derecho sobre los resultados producidos.

%%%%%%%%%%%%%%%%%%%%%%%%%%%%%%%%%%%%%%%%%%%%%%%%%%%%%%%%%%%
\subsection{Consideraciones éticas}
\label{subsec:etica}

El uso de LLM para la generación automática de modelos de procesos de negocio y para el análisis del valor percibido por el cliente conlleva una serie de riesgos éticos que deben tenerse en cuenta en el diseño del sistema:

\begin{itemize}
    \item \textbf{Fiabilidad y supervisión humana}: los LLM pueden generar modelos BPMN incorrectos, incompletos o que no representen fielmente la intención del usuario (fenómeno conocido como \textit{alucinación}). Por ello, el sistema se diseña como una \textbf{herramienta de apoyo} que propone modelos editables por el usuario a través del propio editor visual de \textit{bpmn.io}, en ningún caso como un sistema de decisión autónoma. El usuario mantiene en todo momento el control final sobre el modelo resultante.
    \item \textbf{Transparencia}: el sistema informa al usuario de qué proveedor de LLM se está utilizando en cada momento, de modo que el usuario es consciente de que las respuestas provienen de un sistema de inteligencia artificial generativa y no de un proceso determinista.
    \item \textbf{Sesgos en el análisis PERVAL}: la clasificación del valor percibido (dimensiones emocional, social, de calidad/rendimiento y de precio, según Sweeney y Soutar~\cite{sweeney2001perval}) realizada por el LLM puede verse afectada por los sesgos presentes en los datos de entrenamiento del modelo. Por este motivo, los resultados del análisis PERVAL se presentan como una \textbf{orientación} para el usuario, acompañados de una justificación textual generada por el propio LLM, de modo que el usuario pueda valorar críticamente dicha clasificación.
    \item \textbf{Uso responsable de recursos}: cada interacción con el sistema implica una o varias llamadas a un proveedor externo de LLM, lo que tiene un coste computacional y económico asociado. El diseño del sistema (validación previa de las entradas, mecanismos de \textit{fallback} entre modelos y reparación local del XML generado, en lugar de reintentos completos) busca minimizar el número de llamadas innecesarias al LLM.
\end{itemize}

%%%%%%%%%%%%%%%%%%%%%%%%%%%%%%%%%%%%%%%%%%%%%%%%%%%%%%%%%%%
\section{Identificación y análisis de soluciones posibles}
\label{sec:solucionesPosibles}

Una vez delimitado el alcance funcional del sistema y analizado su marco legal y ético, en esta sección se identifican y comparan las principales alternativas tecnológicas consideradas para abordar el problema planteado. Concretamente, se analizan cuatro decisiones: la selección del editor BPMN sobre el que construir la extensión, la selección del proveedor o proveedores de LLM, la arquitectura de integración entre la extensión y los LLM, y las tecnologías empleadas en el desarrollo del \textit{frontend} y del servidor.

%%%%%%%%%%%%%%%%%%%%%%%%%%%%%%%%%%%%%%%%%%%%%%%%%%%%%%%%%%%
\subsection{Selección del editor BPMN base}
\label{subsec:seleccionEditor}

El primer aspecto a decidir es sobre qué herramienta de modelado BPMN construir la extensión que permita la generación y edición de modelos mediante lenguaje natural. Se consideraron las siguientes alternativas:

\begin{itemize}
    \item \textbf{Desarrollo de un \textit{plugin} para \textit{draw.io}}: \textit{draw.io} (también conocido como \textit{diagrams.net}, integrado en \textit{Google Drive}) es una herramienta de diagramación ampliamente utilizada que admite la edición de diagramas BPMN. Sin embargo, tras estudiar sus mecanismos de extensión se descartó esta alternativa, ya que la plataforma no permite incorporar \textit{plugins} propios desarrollados por terceros: las únicas extensiones disponibles son las ofrecidas oficialmente por la propia herramienta, sin posibilidad de añadir un panel conversacional ni de manipular el modelo mediante código propio.
    \item \textbf{\textit{bpmn.io} (\textit{bpmn-js})}: conjunto de librerías \textit{JavaScript} de código abierto, mantenidas por \textit{Camunda}, que implementan un editor BPMN 2.0 completo ejecutable en el navegador, junto con la librería de bajo nivel \textit{diagram-js} sobre la que se construye. Ofrece un repositorio de ejemplos (\texttt{bpmn-js-examples}) que sirven como punto de partida para extensiones personalizadas, una API completa para manipular el modelo de forma programática (importar/exportar XML, añadir y modificar elementos, gestionar el \textit{canvas}) y un sistema de módulos y extensiones (\textit{additional modules}) que permite añadir funcionalidades sin modificar el núcleo de la librería.
\end{itemize}

Se optó por \textbf{\textit{bpmn-js}}, concretamente partiendo del ejemplo \texttt{modeler} del repositorio \texttt{bpmn-js-examples}, al ser una herramienta de código abierto que sí permite extender el editor con un panel conversacional propio y manipular el modelo BPMN mediante código, por lo que no se plantearon más alternativas. Además, al ser una librería web basada en \textit{JavaScript}, resulta directamente compatible con el resto de tecnologías \textit{frontend} empleadas y con el despliegue de la extensión como sitio estático.

%%%%%%%%%%%%%%%%%%%%%%%%%%%%%%%%%%%%%%%%%%%%%%%%%%%%%%%%%%%
\subsection{Selección del proveedor de LLM}
\label{subsec:seleccionLLM}

La generación de modelos BPMN a partir de lenguaje natural y el análisis PERVAL requieren el uso de un LLM capaz de seguir instrucciones complejas y devolver respuestas estructuradas (XML BPMN o JSON). Se consideraron las siguientes alternativas:

\begin{itemize}
    \item \textbf{Acceso directo a la API de \textit{OpenAI} (\textit{ChatGPT}) o de \textit{Google} (\textit{Gemini})}: ambos proveedores ofrecen modelos de gran capacidad, adecuados tanto para la generación de XML BPMN como para el análisis PERVAL, como muestran trabajos previos sobre generación de modelos BPMN con LLM~\cite{kourani2025evaluating,bise2025bpmn}. Sin embargo, el acceso a sus respectivas API mediante clave propia conlleva un coste económico asociado al volumen de peticiones realizadas, lo que resulta poco adecuado para un proyecto académico que requiere un número elevado de llamadas durante el desarrollo y las pruebas.
    \item \textbf{\textit{GitHub Models}}: servicio de Microsoft/GitHub que ofrece acceso gratuito (dentro de unos límites de uso) a una amplia variedad de modelos de distintos proveedores (\texttt{gpt-4.1}, \texttt{gpt-4o}, \texttt{gpt-4o-mini}, \texttt{gpt-4.1-mini}, \texttt{gpt-4.1-nano}, \texttt{Phi-4}, entre otros) a través de una API compatible con el formato de \textit{OpenAI}.
\end{itemize}

Se optó por \textbf{\textit{GitHub Models}} como proveedor principal de LLM, al ofrecer un plan gratuito que permite realizar un número elevado de peticiones durante el desarrollo y las pruebas sin coste económico, además de dar acceso a varios modelos de distintos proveedores a través de una única API compatible con el formato de \textit{OpenAI}. Dentro de los modelos disponibles, se seleccionó \textbf{\texttt{gpt-4.1}} como modelo principal, por ser el más capaz de la familia GPT disponible en \textit{GitHub Models} en el momento del desarrollo. Para mejorar la resiliencia del sistema, se implementa además un mecanismo de \textbf{selección automática de modelo de respaldo} (\textit{fallback}) que, ante errores o límites de uso del modelo activo, prueba secuencialmente con los siguientes modelos de la lista (\texttt{gpt-4o}, \texttt{gpt-4o-mini}, \texttt{gpt-4.1-mini}, \texttt{gpt-4.1-nano} y \texttt{Phi-4}), de más a menos capaz, lanzando un error únicamente si fallan todos.


%%%%%%%%%%%%%%%%%%%%%%%%%%%%%%%%%%%%%%%%%%%%%%%%%%%%%%%%%%%
\subsection{Arquitectura de integración: extensión directa frente a servidor MCP}
\label{subsec:arquitecturaIntegracion}

Otra decisión relevante es cómo conectar la extensión sobre \textit{bpmn.io}, que se ejecuta en el navegador del usuario, con los proveedores de LLM. Se consideraron dos alternativas principales:

\begin{itemize}
    \item \textbf{Llamadas directas desde el cliente al proveedor de LLM}: la extensión, ejecutándose enteramente en el navegador, invocaría directamente la API del proveedor de LLM. Esta alternativa es la más sencilla de implementar, pero presenta un problema crítico de seguridad: para realizar dichas llamadas sería necesario incluir las claves de API (\textit{tokens}) directamente en el código \textit{JavaScript} del cliente, quedando expuestas a cualquier usuario que inspeccione el código fuente de la página. Además, dificultaría la incorporación de lógica adicional (validación, reparación de XML, selección de modelo de respaldo) sin duplicarla en el cliente.
    \item \textbf{Servidor intermedio mediante el protocolo MCP (\textit{Model Context Protocol})}: se introduce un servidor \textit{backend} (Node.js/Express) que expone un conjunto de herramientas (\textit{tools}) siguiendo la especificación de \textit{Model Context Protocol}~\cite{anthropic2024mcp}, encargado de recibir las peticiones de la extensión, construir los \textit{prompts}, invocar al LLM correspondiente (con su clave de API almacenada de forma segura como variable de entorno del servidor) y devolver el resultado procesado (XML BPMN validado y, en su caso, reparado, o el resultado del análisis PERVAL).
\end{itemize}

Se optó por la segunda alternativa: un \textbf{servidor MCP} como intermediario entre la extensión cliente y los proveedores de LLM. Esta decisión se justifica por varios motivos:

\begin{enumerate}
    \item \textbf{Seguridad}: las claves de API de los proveedores de LLM (\texttt{GITHUB\_TOKEN}) se almacenan exclusivamente como variables de entorno del servidor, sin exponerse en ningún momento al cliente.
    \item \textbf{Centralización de la lógica}: la construcción de \textit{prompts}, el \textit{parsing} robusto de las respuestas del LLM , la reparación de XML BPMN truncado o malformado  y el mecanismo de selección de modelo de respaldo (con \textit{fallback} entre los modelos de \texttt{LLM\_MODELS}) se implementan una única vez en el servidor, en lugar de duplicarse en el cliente.
    \item \textbf{Extensibilidad}: el protocolo MCP define un formato estandarizado de herramientas (\textit{tools}) con esquemas de entrada y salida validados mediante \textit{Zod}, lo que facilita añadir nuevas funcionalidades --como el módulo de análisis PERVAL como nuevas herramientas del mismo servidor, sin modificar la arquitectura general.
    \item \textbf{Despliegue independiente}: al tratarse de dos componentes independientes (extensión estática y servidor MCP), pueden desplegarse por separado en este caso, como dos servicios de \textit{Render} (\texttt{tfm-bpmn-frontend} y \texttt{tfm-bpmn-mcp-server}) y escalarse o sustituirse de forma independiente.
\end{enumerate}

%%%%%%%%%%%%%%%%%%%%%%%%%%%%%%%%%%%%%%%%%%%%%%%%%%%%%%%%%%%
\subsection{Tecnologías para el desarrollo del \textit{frontend} y del servidor}
\label{subsec:tecnologiasFrontend}

Por último, se analizaron las tecnologías concretas a emplear tanto en la extensión cliente como en el servidor MCP:

\begin{itemize}
    \item \textbf{Empaquetado del cliente}: dado que el ejemplo \texttt{modeler} de \textit{bpmn-js-examples} ya utiliza \textit{Webpack} como herramienta de empaquetado, se mantuvo esta elección para minimizar los cambios necesarios sobre la estructura del proyecto y aprovechar la configuración existente para la carga de estilos, iconos y módulos de \textit{bpmn-js}.
    \item \textbf{Interfaz del panel conversacional}: se valoró el uso de un \textit{framework} de componentes (p.~ej. \textit{React} o \textit{Vue}) frente a la integración del panel directamente con \textit{JavaScript} y \textit{jQuery} sobre el DOM existente. Dado que el ejemplo base de \textit{bpmn-js-examples} no utiliza ningún \textit{framework} de este tipo, y que el panel conversacional es un componente relativamente sencillo (lista de mensajes, campo de entrada, indicadores de estado), se optó por \textbf{\textit{JavaScript} y \textit{jQuery}}, evitando así introducir una dependencia adicional de gran tamaño y manteniendo la coherencia con el resto del ejemplo.
    \item \textbf{Servidor MCP}: se eligió \textbf{Node.js con \textit{Express}} para implementar el servidor MCP, dado que es el entorno de ejecución natural para el \textit{SDK} oficial de \textit{Model Context Protocol}~\cite{anthropic2024mcp}, y porque permite compartir el mismo lenguaje (\textit{JavaScript}) entre cliente y servidor. La validación de los esquemas de entrada y salida de cada herramienta MCP se realiza mediante \textbf{\textit{Zod}}, que se integra de forma nativa con el \textit{SDK} de MCP.
    \item \textbf{Internacionalización}: para dar soporte a usuarios de habla española e inglesa, la interfaz y los \textit{prompts} enviados al LLM se parametrizan según el idioma seleccionado (\texttt{es}/\texttt{en}), de modo que tanto los textos de la interfaz como las respuestas del LLM (incluidas las justificaciones del análisis PERVAL) se generen en el idioma elegido por el usuario.
    \item \textbf{Pruebas}: se incorporó \textbf{\textit{Playwright}} para la realización de pruebas \textit{end-to-end} sobre la extensión, permitiendo automatizar la verificación de los flujos principales (generación de un modelo a partir de una descripción, edición conversacional, análisis PERVAL) en un navegador real.
    \item \textbf{Despliegue}: se seleccionó \textbf{\textit{Render}} como plataforma de despliegue, por ofrecer un nivel gratuito adecuado para un proyecto académico y por permitir definir, mediante un único archivo \texttt{render.yaml}, los dos servicios necesarios: un sitio estático para la extensión (\texttt{tfm-bpmn-frontend}) y un servicio \textit{web} de tipo Node.js para el servidor MCP (\texttt{tfm-bpmn-mcp-server}).
\end{itemize}

{Como resultado del análisis realizado en las secciones anteriores, se seleccionaron las alternativas tecnológicas que mejor satisfacen los requisitos funcionales y no funcionales identificados para el desarrollo de la solución propuesta. En todos los casos, la elección se fundamentó en criterios relacionados con la interoperabilidad, la extensibilidad, la mantenibilidad y la adecuación a los objetivos del presente TFM. La Tabla \ref{tab:design_decisions} resume las principales decisiones de diseño adoptadas, las alternativas consideradas durante el análisis y los criterios que justificaron cada elección. Sobre esta base, en la siguiente sección se presenta la solución propuesta.}


\begin{table}[htbp]
\centering
\caption{Resumen de las principales decisiones de diseño adoptadas.}
\label{tab:design_decisions}
\renewcommand{\arraystretch}{1.25}
\footnotesize
\begin{tabularx}{\textwidth}{|p{3.1cm}|X|X|p{3.1cm}|}
\hline
\textbf{Decisión de diseño} &
\textbf{Alternativas consideradas} &
\textbf{Criterio principal de evaluación} &
\textbf{Decisión adoptada} \\
\hline

Editor de modelado BPMN &
draw.io, Camunda Modeler, bpmn.io &
Capacidad de extensión, edición interactiva, acceso al modelo BPMN y disponibilidad de API. &
\textbf{bpmn.io} \\
\hline

Proveedor de LLM &
API de OpenAI (ChatGPT), API de Google (Gemini), GitHub Models &
Disponibilidad de plan gratuito, acceso a múltiples modelos y API unificada compatible con el formato OpenAI. &
\textbf{GitHub Models} (modelo principal \texttt{gpt-4.1} con mecanismo de \textit{fallback} automático) \\
\hline

Arquitectura de integración &
Llamadas directas del cliente al proveedor de LLM; servidor MCP intermedio (Node.js/Express) &
Seguridad de credenciales, centralización de la lógica de construcción de \textit{prompts} y reparación de XML. &
\textbf{Servidor MCP} (Node.js + Express) \\
\hline

Comunicación con el LLM &
SDK nativo de cada proveedor (OpenAI SDK, Google AI SDK\ldots); API REST compatible con formato OpenAI &
Unificación de llamadas multi-modelo y compatibilidad directa con GitHub Models sin dependencias adicionales. &
\textbf{API REST compatible con OpenAI} (punto de acceso único para todos los modelos de GitHub Models) \\
\hline

Tecnologías de desarrollo &
React/Vue + Python/Flask; JavaScript + Node.js/Express; otras combinaciones &
Coherencia con el ejemplo base de \textit{bpmn-js}, lenguaje compartido cliente-servidor e integración nativa con el MCP SDK. &
\textbf{JavaScript}: cliente (bpmn-js, jQuery, Webpack); servidor (Node.js, Express, Zod) \\
\hline

\end{tabularx}
\end{table}

\denis{Las decisiones sintetizadas en la Tabla \ref{tab:design_decisions} constituyen la base tecnológica sobre la que se construye la solución desarrollada en este trabajo. A partir de ellas se derivan los requisitos funcionales y no funcionales que se presentan en la siguiente sección, los cuales concretan el comportamiento esperado del sistema y las características de calidad que deberá satisfacer.}

%%%%%%%%%%%%%%%%%%%%%%%%%%%%%%%%%%%%%%%%%%%%%%%%%%%%%%%%%%%
\section{Solución propuesta}
\label{sec:solucionPropuesta}

A partir del modelado funcional y de las decisiones tecnológicas descritas anteriormente, se concreta la solución propuesta mediante la especificación de los requisitos funcionales (tabla~\ref{tab:requisitosFuncionales}) y no funcionales (tabla~\ref{tab:requisitosNoFuncionales}) del sistema, y mediante un conjunto de quince historias de usuario (US01-US15) que detallan el trabajo realizado para satisfacer dichos requisitos.

%%%%%%%%%%%%%%%%%%%%%%%%%%%%%%%%%%%%%%%%%%%%%%%%%%%%%%%%%%%
\subsection{Requisitos funcionales}
\label{subsec:requisitosFuncionales}

La tabla~\ref{tab:requisitosFuncionales} recoge los requisitos funcionales (RF) identificados para el sistema.

\begin{longtable}{|c|p{11.5cm}|}
\hline
\rowcolor{gray!20} \textbf{Código} & \textbf{Descripción} \\
\hline
\endfirsthead
\hline
\rowcolor{gray!20} \textbf{Código} & \textbf{Descripción} \\
\hline
\endhead
RF-01 & El sistema debe permitir al usuario describir un proceso de negocio en lenguaje natural y generar, a partir de dicha descripción, un modelo BPMN 2.0 válido. \\
\hline
RF-02 & El sistema debe permitir al usuario solicitar modificaciones sobre un modelo BPMN existente mediante instrucciones en lenguaje natural, conservando el contexto de la conversación. \\
\hline
RF-03 & El sistema debe visualizar el modelo BPMN generado o modificado en el lienzo de \textit{bpmn.io}. \\
\hline
RF-04 & El sistema debe permitir exportar el modelo BPMN resultante en formato XML. \\
\hline
RF-05 & El sistema debe validar y, en caso de errores menores, reparar automáticamente el XML BPMN devuelto por el LLM antes de cargarlo en el editor. \\
\hline
RF-06 & El sistema debe seleccionar automáticamente un modelo de respaldo cuando el modelo principal no esté disponible o devuelva un error. \\
\hline
RF-07 & El sistema debe permitir analizar el valor percibido (PERVAL) de las tareas o entregables de un modelo BPMN, según las dimensiones emocional, social, de calidad/rendimiento y de precio~\cite{sweeney2001perval}. \\
\hline
RF-08 & El sistema debe permitir seleccionar el ámbito (tareas o entregables) y el actor sobre el que se realiza el análisis PERVAL. \\
\hline
RF-09 & El sistema debe presentar visualmente los resultados del análisis PERVAL (resumen por dimensión y valor general), acompañados de una justificación textual. \\
\hline
RF-10 & El sistema debe permitir seleccionar el idioma de la interfaz y de las respuestas generadas (español/inglés). \\
\hline
RF-11 & El sistema debe informar al usuario del estado de cada petición (cargando, completada, error) en todo momento. \\
\hline
\caption{Requisitos funcionales del sistema.}
\label{tab:requisitosFuncionales}
\end{longtable}


%%%%%%%%%%%%%%%%%%%%%%%%%%%%%%%%%%%%%%%%%%%%%%%%%%%%%%%%%%%
\subsection{Requisitos no funcionales}
\label{subsec:requisitosNoFuncionales}

La tabla~\ref{tab:requisitosNoFuncionales} recoge los requisitos no funcionales (RNF) identificados para el sistema.

\begin{longtable}{|c|p{11.5cm}|}
\hline
\rowcolor{gray!20} \textbf{Código} & \textbf{Descripción} \\
\hline
\endfirsthead
\hline
\rowcolor{gray!20} \textbf{Código} & \textbf{Descripción} \\
\hline
\endhead
RNF-01 & \textbf{Usabilidad}: la interacción con el sistema debe realizarse mediante un panel conversacional integrado en el propio editor de \textit{bpmn.io}, sin requerir conocimientos previos de la notación BPMN ni de los proveedores de LLM. \\
\hline
RNF-02 & \textbf{Seguridad}: las claves de API de los proveedores de LLM deben almacenarse únicamente como variables de entorno del servidor MCP, sin exponerse en ningún momento al cliente. \\
\hline
RNF-03 & \textbf{Tiempo de respuesta}: el sistema debe proporcionar realimentación visual inmediata (indicadores de carga) ante cualquier petición al LLM, dado que los tiempos de respuesta de estos servicios pueden oscilar entre varios segundos y, en modelos complejos, más de un minuto. \\
\hline
RNF-04 & \textbf{Extensibilidad}: el servidor MCP debe organizarse como un conjunto de herramientas (\textit{tools}) independientes con esquemas de entrada/salida validados mediante \textit{Zod}, de modo que puedan añadirse nuevas funcionalidades sin modificar la arquitectura general. \\
\hline
RNF-05 & \textbf{Portabilidad}: el sistema debe poder desplegarse en una plataforma \textit{cloud} (\textit{Render}) mediante un único fichero de configuración (\texttt{render.yaml}), y ser ejecutable en cualquier entorno compatible con Node.js. \\
\hline
RNF-06 & \textbf{Mantenibilidad}: el código debe organizarse de forma modular, separando claramente la extensión cliente, el servidor MCP, la lógica de llamada a los LLM (\texttt{llm.js}) y cada herramienta MCP (\texttt{tools/}). \\
\hline

RNF-07 & \textbf{Conformidad}: los modelos BPMN generados o modificados deben generarse en un formato compatible con BPMN 2.0~\cite{omg2011bpmn} y susceptible de ser importado por herramientas compatibles. \\
\hline
RNF-08 & \textbf{Privacidad}: el sistema debe minimizar los datos enviados a los proveedores de LLM (principio de minimización de datos del RGPD~\cite{rgpd2016}) y no debe almacenar de forma persistente datos personales de los usuarios. \\
\hline
\caption{Requisitos no funcionales del sistema.}
\label{tab:requisitosNoFuncionales}
\end{longtable}


%%%%%%%%%%%%%%%%%%%%%%%%%%%%%%%%%%%%%%%%%%%%%%%%%%%%%%%%%%%
\subsection{Priorización de requisitos. MoSCoW}
\label{subsec:priorizacion}
%%%%%%%%%%%%%%%%%%%%%%%%%%%%%%%%%%%%%%%%%%%%%%%%%%%%%%%%%%%

Con el objetivo de establecer un orden claro de implementación y gestionar los recursos disponibles dentro de los límites de un proyecto académico, se ha aplicado la metodología MoSCoW para priorizar los requisitos identificados en las secciones anteriores. Esta metodología clasifica los requisitos en cuatro categorías según su importancia crítica para el éxito del proyecto:

\begin{itemize}
    \item[$\bullet$] \textbf{Must have} (debe tener): funcionalidades absolutamente esenciales cuya ausencia impediría que el sistema cumpla su propósito fundamental.
    \item[$\bullet$] \textbf{Should have} (debería tener): funcionalidades importantes que aportan valor significativo pero cuya ausencia no impide el funcionamiento básico del sistema.
    \item[$\bullet$] \textbf{Could have} (podría tener): funcionalidades deseables que pueden implementarse si los recursos lo permiten, sin comprometer los objetivos principales.
    \item[$\bullet$] \textbf{Won't have} (no tendrá): funcionalidades descartadas explícitamente para la versión actual del proyecto, ya sea por limitaciones de tiempo, recursos o decisiones estratégicas de desarrollo.
\end{itemize}

\subsubsection*{Must have}

Se consideran imprescindibles aquellos requisitos sin los cuales el sistema no puede cumplir su propósito fundamental: la generación y el análisis iterativo de modelos BPMN mediante LLM.

La generación de modelos BPMN a partir de descripciones en lenguaje natural (RF-01) constituye el núcleo de la propuesta; sin esta funcionalidad, la herramienta pierde su razón de ser. Íntimamente ligada a ella, la modificación iterativa del modelo mediante instrucciones conversacionales (RF-02) es igualmente esencial, ya que materializa el paradigma de \textit{vibe modeling} sobre el que se fundamenta el proyecto e impide que el sistema sea un mero generador estático de diagramas. La visualización del modelo resultante en el lienzo de \textit{bpmn.io} (RF-03) es también imprescindible: sin ella el usuario no puede validar el resultado generado ni interactuar con el editor.

La validación y reparación automática del XML BPMN (RF-05) es un requisito crítico, dado que los LLMs producen con frecuencia XML malformado o truncado; sin este mecanismo, los modelos generados serían inutilizables en la mayoría de los casos. El análisis del valor percibido según la clasificación PERVAL (RF-07) constituye el segundo objetivo central del proyecto académico; su ausencia eliminaría la componente de análisis orientado al cliente que diferencia esta herramienta de otros asistentes de modelado basados en LLM. La presentación visual de los resultados del análisis con justificación textual (RF-09) es igualmente esencial: sin ella el análisis PERVAL no aporta valor interpretable al usuario.

En cuanto a los requisitos no funcionales, la integración del panel conversacional directamente en el entorno de \textit{bpmn.io} sin requerir conocimientos previos de la notación BPMN (RNF-01) es el requisito de usabilidad fundamental que materializa el concepto de \textit{vibe modeling}. El almacenamiento de las claves de API exclusivamente como variables de entorno del servidor MCP (RNF-02) es crítico desde el punto de vista de la seguridad: exponer dichas credenciales en el código del cliente pondría en riesgo las cuentas de los usuarios y comprometería la viabilidad de la herramienta en producción.

\subsubsection*{Should have}

Esta categoría engloba requisitos que aportan valor significativo a la utilidad, la resiliencia y la calidad del sistema, sin ser estrictamente imprescindibles para el funcionamiento básico.

La exportación del modelo BPMN en formato XML (RF-04) incrementa considerablemente la utilidad práctica de la herramienta, al permitir al usuario conservar y reutilizar los modelos fuera del entorno de \textit{bpmn.io}. El mecanismo de selección automática de modelo de respaldo (RF-06) es importante para garantizar la continuidad del servicio ante los límites de uso de la API de \textit{GitHub Models}, especialmente frecuentes durante el desarrollo y las pruebas; sin él, el sistema fallaría con mayor frecuencia ante picos de peticiones. La configuración del ámbito (tareas o entregables) y del actor sobre el que se realiza el análisis PERVAL (RF-08) es relevante para adaptar el análisis a distintos contextos de proceso, aunque podría implementarse con valores predeterminados en una primera versión. La información del estado de cada petición mediante indicadores visuales (RF-11) mejora notablemente la experiencia de usuario en interacciones con latencias elevadas, habituales en llamadas a LLMs.

Desde el punto de vista técnico, la organización del servidor MCP como conjunto de herramientas independientes con esquemas validados mediante \textit{Zod} (RNF-04) facilita la incorporación de nuevas funcionalidades sin alterar la arquitectura general. La modularidad del código, con separación clara entre el cliente, el servidor MCP, el módulo de llamada al LLM y las herramientas individuales (RNF-06), facilita el mantenimiento y la evolución del sistema. La conformidad de los modelos generados con el estándar BPMN 2.0 (RNF-07) garantiza la interoperabilidad con otras herramientas de modelado, aunque está parcialmente cubierta por el mecanismo de validación de RF-05.

\subsubsection*{Could have}

Los siguientes requisitos añaden valor pero pueden implementarse de forma diferida sin comprometer los objetivos principales del proyecto.

La selección del idioma de la interfaz y de las respuestas generadas (RF-10) amplía el alcance de la herramienta a usuarios de diferentes entornos lingüísticos. Sin embargo, el sistema es plenamente funcional en un único idioma, y su implementación puede postergarse hasta que las funcionalidades principales estén consolidadas. La realimentación visual inmediata ante cualquier petición al LLM —más allá del indicador de carga básico— (RNF-03) mejoraría la percepción de la capacidad de respuesta del sistema, pero no afecta a su comportamiento funcional. La portabilidad del sistema mediante un único fichero de configuración para su despliegue en \textit{Render} (RNF-05) es deseable para las pruebas de aceptación TAM con usuarios reales, aunque durante el desarrollo basta con ejecutar el sistema en local. La gestión explícita de la minimización de datos personales conforme al RGPD (RNF-08) está parcialmente cubierta por el requisito de seguridad de las claves de API (RNF-02); los mecanismos adicionales de privacidad son deseables pero no bloquean el funcionamiento de la herramienta.

\subsubsection*{Won't have}

Esta categoría incluye funcionalidades descartadas de forma explícita para la versión actual del proyecto, ya sea por limitaciones de tiempo y recursos, o por decisiones estratégicas orientadas a validar el concepto fundamental antes de ampliar el alcance.

El soporte para proveedores de LLM adicionales distintos de \textit{GitHub Models} fue considerado durante el análisis de alternativas y descartado (véase \S\ref{subsec:seleccionLLM}): el plan gratuito de \textit{GitHub Models} ofrece acceso a modelos GPT de alta calidad con límites suficientes para los objetivos del proyecto, y añadir proveedores adicionales incrementaría la complejidad del servidor MCP sin aportar valor diferencial en esta primera versión. La persistencia del historial de conversación entre sesiones fue igualmente descartada: el historial se gestiona en memoria del cliente durante la sesión activa, lo cual es suficiente para los ciclos de refinamiento iterativo previstos; una persistencia duradera requeriría una base de datos, lo que excede el alcance del proyecto. Por último, la exportación del diagrama en formatos alternativos (imagen PNG/SVG, PDF) no forma parte de los objetivos del TFM y puede cubrirse mediante las funcionalidades nativas del propio \textit{bpmn.io}.

\vspace{0.5em}
\noindent\textbf{Resumen por prioridad:}

Must (8): RF-01, RF-02, RF-03, RF-05, RF-07, RF-09, RNF-01, RNF-02

Should (7): RF-04, RF-06, RF-08, RF-11, RNF-04, RNF-06, RNF-07

Could (4): RF-10, RNF-03, RNF-05, RNF-08

Won't (3 funcionalidades): integración con otros proveedores de LLM, persistencia del historial entre sesiones, exportación en formatos alternativos

\vspace{0.5em}
\begingroup\small
\begin{longtable}{|p{4.2cm}|p{2.3cm}|p{2.3cm}|p{2.3cm}|p{2.3cm}|}
\hline
\rowcolor{gray!20} \textbf{Funcionales} & \textbf{Must} & \textbf{Should} & \textbf{Could} & \textbf{Won't} \\
\hline
\endfirsthead
\hline
\rowcolor{gray!20} \textbf{Funcionales} & \textbf{Must} & \textbf{Should} & \textbf{Could} & \textbf{Won't} \\
\hline
\endhead
Generación BPMN (RF-01 a RF-06) & 01, 02, 03, 05 & 04, 06 & --- & --- \\
\hline
Análisis PERVAL (RF-07 a RF-09) & 07, 09 & 08 & --- & --- \\
\hline
Configuración/UX (RF-10 a RF-11) & --- & 11 & 10 & --- \\
\hline
\caption{Resumen de priorización MoSCoW de requisitos funcionales.}
\label{tab:moscowFuncionales}
\end{longtable}
\endgroup

\begingroup\small
\begin{longtable}{|p{4.2cm}|p{2.3cm}|p{2.3cm}|p{2.3cm}|p{2.3cm}|}
\hline
\rowcolor{gray!20} \textbf{No funcionales} & \textbf{Must} & \textbf{Should} & \textbf{Could} & \textbf{Won't} \\
\hline
\endfirsthead
\hline
\rowcolor{gray!20} \textbf{No funcionales} & \textbf{Must} & \textbf{Should} & \textbf{Could} & \textbf{Won't} \\
\hline
\endhead
Usabilidad (RNF-01) & 01 & --- & --- & --- \\
\hline
Seguridad y privacidad (RNF-02, 08) & 02 & --- & 08 & --- \\
\hline
Rendimiento (RNF-03) & --- & --- & 03 & --- \\
\hline
Arquitectura y calidad (RNF-04 a RNF-07) & --- & 04, 06, 07 & 05 & --- \\
\hline
\caption{Resumen de priorización MoSCoW de requisitos no funcionales.}
\label{tab:moscowNoFuncionales}
\end{longtable}
\endgroup

%%%%%%%%%%%%%%%%%%%%%%%%%%%%%%%%%%%%%%%%%%%%%%%%%%%%%%%%%%%
\subsection{Historias de usuario}
\label{subsec:historiasUsuario}

Para concretar el desarrollo del sistema se ha definido un conjunto de quince historias de usuario (US01-US15), siguiendo una metodología ágil basada en \textit{Scrum}~\cite{riano2021}. Para cada historia de usuario se indica una breve descripción, una estimación del esfuerzo necesario (en horas) y el desglose en tareas concretas, también estimadas en horas. Estas historias de usuario constituyen la base de la planificación del trabajo por \textit{sprints} presentada a continuación.

\textbf{US01 -- Selección tecnológica y planificación del proyecto}

\textit{Descripción}: como desarrollador, quiero analizar las alternativas tecnológicas disponibles para la generación de modelos BPMN mediante LLM y para el análisis PERVAL, para seleccionar el editor BPMN base, el o los proveedores de LLM, y la arquitectura de integración, y planificar el desarrollo del proyecto por \textit{sprints}.

\textit{Estimación}: 8h.

\textit{Tareas}:
\begin{itemize}
    \item Estudio de trabajos relacionados sobre generación de modelos BPMN mediante LLM (3h).
    \item Selección del editor BPMN, del proveedor de LLM y de la arquitectura de integración (3h).
    \item Definición del \textit{backlog} inicial y planificación de \textit{sprints} (2h).
\end{itemize}

\textbf{US02 -- Configuración del entorno de desarrollo}

\textit{Descripción}: como desarrollador, quiero configurar un entorno de desarrollo a partir del ejemplo \texttt{modeler} de \texttt{bpmn-js-examples}, para que disponga de un editor BPMN funcional sobre el que construir la extensión.

\textit{Estimación}: 6h.

\textit{Tareas}:
\begin{itemize}
    \item Clonado y configuración inicial del proyecto \texttt{modeler} (2h).
    \item Configuración de \textit{Webpack} y de las dependencias del proyecto (2h).
    \item Verificación del funcionamiento del editor BPMN base (2h).
\end{itemize}

\textbf{US03 -- Diseño e implementación del servidor MCP}

\textit{Descripción}: como desarrollador, quiero implementar un servidor \textit{backend} que siga la especificación de \textit{Model Context Protocol}~\cite{anthropic2024mcp}, para que actúe como intermediario seguro entre la extensión cliente y los proveedores de LLM.

\textit{Estimación}: 12h.

\textit{Tareas}:
\begin{itemize}
    \item Estudio de la especificación MCP y del \textit{SDK} oficial (3h).
    \item Implementación del servidor Express con transporte HTTP y registro de herramientas (4h).
    \item Definición de la estructura de herramientas (\textit{tools}) y de sus esquemas de entrada/salida con \textit{Zod} (3h).
    \item Pruebas de conexión entre la extensión cliente y el servidor MCP (2h).
\end{itemize}

\textbf{US04 -- Integración con \textit{GitHub Models}}

\textit{Descripción}: como desarrollador, quiero integrar el servidor MCP con la API de \textit{GitHub Models}, para poder invocar distintos modelos de LLM (\texttt{gpt-4.1}, \texttt{gpt-4o}, etc.) para la generación de modelos BPMN y el análisis PERVAL.

\textit{Estimación}: 10h.

\textit{Tareas}:
\begin{itemize}
    \item Configuración de credenciales (\texttt{GITHUB\_TOKEN}) y definición del catálogo de modelos disponibles (2h).
    \item Implementación de la función de llamada al LLM (\texttt{callOpenAI}) (4h).
    \item Implementación del mecanismo de selección de modelo de respaldo ante errores (2h).
    \item Pruebas de integración con distintos modelos (2h).
\end{itemize}

\textbf{US05 -- Generación de modelos BPMN a partir de lenguaje natural}

\textit{Descripción}: como usuario, quiero describir un proceso de negocio en lenguaje natural y obtener automáticamente un modelo BPMN 2.0 que lo represente, para poder partir de un primer borrador del modelo sin tener que construirlo manualmente.

\textit{Estimación}: 14h.

\textit{Tareas}:
\begin{itemize}
    \item Diseño del \textit{prompt} de generación, incluyendo el catálogo completo de elementos BPMN soportados (4h).
    \item Implementación de la herramienta MCP de generación de modelos BPMN (4h).
    \item Implementación del \textit{parsing} robusto de la respuesta del LLM (\texttt{parseLLMJson}) (2h).
    \item Carga del XML generado en el editor \textit{bpmn-js} (2h).
    \item Pruebas con descripciones de proceso de distinta complejidad (2h).
\end{itemize}

\textbf{US06 -- Edición conversacional iterativa del modelo}

\textit{Descripción}: como usuario, quiero poder solicitar modificaciones sobre el modelo BPMN actual mediante mensajes en lenguaje natural, para poder refinar iterativamente el modelo a través de una conversación con el sistema.

\textit{Estimación}: 12h.

\textit{Tareas}:
\begin{itemize}
    \item Diseño del panel conversacional (historial de mensajes y campo de entrada) (3h).
    \item Implementación de la herramienta MCP de edición, que recibe el XML actual y la instrucción del usuario (4h).
    \item Gestión del historial de la conversación en el cliente (3h).
    \item Pruebas de iteraciones sucesivas de modificación sobre un mismo modelo (2h).
\end{itemize}

\textbf{US07 -- Validación y reparación automática del XML BPMN}

\textit{Descripción}: como usuario, quiero que el sistema detecte y corrija automáticamente los errores habituales en el XML BPMN devuelto por el LLM (etiquetas sin cerrar, XML truncado), para que el modelo se cargue correctamente en el editor incluso cuando la respuesta del LLM no es perfecta.

\textit{Estimación}: 10h.

\textit{Tareas}:
\begin{itemize}
    \item Análisis de los errores más habituales en las respuestas del LLM (truncamientos, comillas sin cerrar, etiquetas incompletas) (3h).
    \item Implementación de la función de limpieza y reparación de XML (\texttt{cleanXML}) (4h).
    \item Implementación de la detección automática de descripciones completas frente a incrementales (2h).
    \item Pruebas con respuestas malformadas del LLM (1h).
\end{itemize}



\textbf{US08 -- Diseño de la herramienta de análisis PERVAL}

\textit{Descripción}: como desarrollador, quiero diseñar una herramienta MCP capaz de analizar el valor percibido por el cliente de las tareas o entregables de un modelo BPMN, según la clasificación de Sweeney y Soutar~\cite{sweeney2001perval}, para que el usuario pueda obtener información sobre el valor que aporta cada elemento del proceso.

\textit{Estimación}: 10h.

\textit{Tareas}:
\begin{itemize}
    \item Estudio de las dimensiones de valor percibido (emocional, social, calidad/rendimiento, precio) (2h).
    \item Diseño del esquema de entrada/salida de la herramienta \texttt{analyzePerval} (3h).
    \item Diseño del \textit{prompt} de análisis PERVAL, incluyendo el ámbito (tareas/entregables) y el actor (3h).
    \item Pruebas con modelos BPMN de ejemplo (2h).
\end{itemize}


\textbf{US09 -- Implementación de la herramienta \texttt{analyzePerval}}

\textit{Descripción}: como desarrollador, quiero implementar en el servidor MCP la herramienta de análisis PERVAL diseñada en US08, para que, dado un modelo BPMN y un conjunto de tareas, el sistema devuelva la clasificación por dimensiones, el resumen agregado y el valor general.

\textit{Estimación}: 10h.

\textit{Tareas}:
\begin{itemize}
    \item Implementación de la herramienta MCP \texttt{analyzePerval} con su esquema \textit{Zod} (3h).
    \item Extracción de las tareas y entregables a partir del XML BPMN (3h).
    \item Cálculo del resumen por dimensión y del valor general a partir de la respuesta del LLM (2h).
    \item Pruebas de la herramienta con distintos ámbitos y actores (2h).
\end{itemize}

\textbf{US10 -- Visualización de los resultados del análisis PERVAL}

\textit{Descripción}: como usuario, quiero visualizar de forma clara los resultados del análisis PERVAL (resumen por dimensión, valor general y justificación de cada tarea), para poder interpretar fácilmente el valor percibido de cada elemento del proceso.

\textit{Estimación}: 10h.

\textit{Tareas}:
\begin{itemize}
    \item Diseño de la interfaz de selección de ámbito (tareas/entregables) y actor (2h).
    \item Implementación de la visualización del resumen por dimensión y del valor general (4h).
    \item Resaltado sobre el diagrama BPMN de los elementos analizados (2h).
    \item Pruebas de usabilidad de la visualización (2h).
\end{itemize}

\textbf{US11 -- Diseño e implementación de la interfaz del panel conversacional}

\textit{Descripción}: como usuario, quiero disponer de un panel conversacional integrado en el editor con un diseño visual coherente y agradable, para que la interacción con el sistema resulte intuitiva y profesional.

\textit{Estimación}: 10h.

\textit{Tareas}:
\begin{itemize}
    \item Diseño de la paleta de colores (azul/verde) y selección de tipografía (\textit{Inter}) (2h).
    \item Implementación de los estilos del panel conversacional, adaptados a distintos tamaños de pantalla (4h).
    \item Implementación de los indicadores de estado de carga y de error (2h).
    \item Pruebas de la interfaz en distintos tamaños de pantalla (2h).
\end{itemize}

\textbf{US12 -- Internacionalización de la interfaz y de los \textit{prompts}}

\textit{Descripción}: como usuario, quiero poder utilizar el sistema en español o en inglés, tanto en la interfaz como en las respuestas generadas por el LLM, para poder trabajar en el idioma de su preferencia.

\textit{Estimación}: 8h.

\textit{Tareas}:
\begin{itemize}
    \item Extracción de los textos de la interfaz a ficheros de idioma (es/en) (3h).
    \item Parametrización de los \textit{prompts} de generación, edición y análisis PERVAL según el idioma seleccionado (3h).
    \item Pruebas del sistema completo en ambos idiomas (2h).
\end{itemize}

\textbf{US13 -- Pruebas \textit{end-to-end} con \textit{Playwright}}

\textit{Descripción}: como desarrollador, quiero disponer de un conjunto de pruebas automatizadas \textit{end-to-end}, para poder verificar que los flujos principales del sistema (generación, edición conversacional y análisis PERVAL) funcionan correctamente en un navegador real.

\textit{Estimación}: 8h.

\textit{Tareas}:
\begin{itemize}
    \item Configuración de \textit{Playwright} sobre el proyecto (2h).
    \item Implementación de pruebas de generación y edición conversacional de modelos BPMN (3h).
    \item Implementación de pruebas de la herramienta \texttt{analyzePerval} (2h).
    \item Integración de las pruebas en el flujo de desarrollo (1h).
\end{itemize}


\textbf{US14 -- Despliegue en \textit{Render} y documentación}

\textit{Descripción}: como desarrollador, quiero desplegar la extensión y el servidor MCP en \textit{Render}, para que el sistema sea accesible públicamente y pueda ser utilizado y evaluado fuera del entorno de desarrollo local.

\textit{Estimación}: 7h.

\textit{Tareas}:
\begin{itemize}
    \item Configuración del fichero \texttt{render.yaml} con los dos servicios (\texttt{tfm-bpmn-frontend} y \texttt{tfm-bpmn-mcp-server}) (2h).
    \item Despliegue y pruebas de funcionamiento del servidor MCP y de la extensión en \textit{Render} (3h).
    \item Documentación de las variables de entorno necesarias y del procedimiento de puesta en marcha (2h).
\end{itemize}

\textbf{US15 -- Validación de la solución con usuarios mediante TAM}

\textit{Descripción}: como desarrollador, quiero validar la solución desarrollada mediante una evaluación con un grupo de aproximadamente quince usuarios, empleando el modelo de aceptación de la tecnología (\textit{Technology Acceptance Model}, TAM), para poder valorar la utilidad percibida, la facilidad de uso percibida y la intención de uso del sistema.

\textit{Estimación}: 10h.

\textit{Tareas}:
\begin{itemize}
    \item Diseño del cuestionario TAM (utilidad percibida, facilidad de uso percibida e intención de uso) adaptado al sistema desarrollado (2h).
    \item Selección y convocatoria de un grupo de aproximadamente quince participantes (2h).
    \item Realización de las sesiones de prueba, en las que los participantes utilizan el sistema desplegado y completan el cuestionario (4h).
    \item Análisis de los resultados del cuestionario y elaboración de conclusiones sobre la aceptación de la solución (2h).
\end{itemize}

La estimación total del conjunto de historias de usuario asciende a \textbf{145 horas}, que se distribuyen en ocho \textit{sprints} semanales tal como se detalla a continuación.

%%%%%%%%%%%%%%%%%%%%%%%%%%%%%%%%%%%%%%%%%%%%%%%%%%%%%%%%%%%
\section{Plan de trabajo}
\label{sec:planTrabajo}

El desarrollo del sistema se ha planificado siguiendo una metodología ágil basada en \textit{Scrum}~\cite{riano2021}, organizando el trabajo en \textbf{ocho \textit{sprints} semanales} (Sprint 0 a Sprint 7), desde el 26/05/2026 hasta el 20/07/2026. Cada \textit{sprint} agrupa una o varias de las historias de usuario definidas en la sección~\ref{subsec:historiasUsuario}, de modo que la suma de las estimaciones de las historias de usuario de cada \textit{sprint} determina su carga de trabajo prevista. De forma paralela al desarrollo de cada \textit{sprint}, se redacta la documentación correspondiente de la memoria del TFM, de modo que, una vez finalizado el Sprint 7, el trabajo restante se centra en completar y revisar las secciones de la memoria que no hayan podido finalizarse durante el desarrollo.

Conviene señalar que la planificación mediante sprints constituye exclusivamente la estrategia adoptada para organizar el desarrollo incremental del artefacto software. La evaluación científica de la propuesta se aborda de forma independiente mediante un estudio empírico de aceptación basado en el TAM, descrito en los capítulos posteriores. De este modo, Scrum se emplea como metodología de desarrollo del software, mientras que la validación de la investigación se fundamenta en un procedimiento empírico específico.

La tabla~\ref{tab:planTrabajo} resume la planificación de los ocho \textit{sprints}, indicando para cada uno su rango de fechas, las historias de usuario asignadas y la carga de trabajo estimada.

\begingroup
\small
\begin{longtable}{|c|c|p{6cm}|c|}
\hline
\rowcolor{gray!20} \textbf{Sprint} & \textbf{Fechas} & \textbf{Historias de usuario} & \textbf{Carga (h)} \\
\hline
\endfirsthead
\hline
\rowcolor{gray!20} \textbf{Sprint} & \textbf{Fechas} & \textbf{Historias de usuario} & \textbf{Carga (h)} \\
\hline
\endhead
Sprint 0 & 26/05/2026 -- 01/06/2026 & US01 (Selección tecnológica y planificación), US02 (Configuración del entorno de desarrollo) & 14 \\
\hline
Sprint 1 & 02/06/2026 -- 08/06/2026 & US03 (Servidor MCP), US04 (Integración con \textit{GitHub Models}) & 22 \\
\hline
Sprint 2 & 09/06/2026 -- 15/06/2026 & US05 (Generación de modelos BPMN a partir de lenguaje natural) & 14 \\
\hline
Sprint 3 & 16/06/2026 -- 22/06/2026 & US06 (Edición conversacional iterativa), US07 (Validación y reparación de XML) & 22 \\
\hline
Sprint 4 & 23/06/2026 -- 29/06/2026 & US08 (Diseño de la herramienta PERVAL), US09 (Implementación de \texttt{analyzePerval}) & 20 \\
\hline
Sprint 5 & 30/06/2026 -- 06/07/2026 & US10 (Visualización de resultados PERVAL), US11 (Interfaz del panel conversacional) & 20 \\
\hline
Sprint 6 & 07/07/2026 -- 13/07/2026 & US12 (Internacionalización), US13 (Pruebas \textit{end-to-end}) & 16 \\
\hline
Sprint 7 & 14/07/2026 -- 20/07/2026 & US14 (Despliegue en \textit{Render} y documentación), US15 (Validación con usuarios mediante TAM) & 17 \\
\hline
\caption{Planificación del trabajo por \textit{sprints}.}
\label{tab:planTrabajo}
\end{longtable}
\endgroup

Como puede observarse, la carga de trabajo prevista se sitúa en torno a las 18 horas semanales de media (145 horas en total repartidas en 8 \textit{sprints}), con dos \textit{sprints} de menor carga (Sprint 0 y Sprint 2) dedicados, respectivamente, a la planificación inicial del proyecto y a la implementación de la funcionalidad central de generación de modelos BPMN, que por su complejidad se ha mantenido como única historia de usuario del \textit{sprint}. Los \textit{sprints} 1, 3, 4 y 5 concentran la mayor parte del desarrollo de las dos funcionalidades principales del sistema --la generación y edición conversacional de modelos BPMN (\textit{sprints} 1-3) y la herramienta de análisis PERVAL (\textit{sprints} 4-5)--, mientras que los \textit{sprints} 6 y 7 se dedican a pulir la interfaz, dar soporte a la internacionalización, completar las pruebas automatizadas, desplegar el sistema en producción y validar la solución con un grupo de usuarios mediante el modelo TAM.

\shorthandon{<>}

%%%%%%%%%%%%%%%%%%%%%%%%%%%%%%%%%%%%%%%%%%%%%%%%%%%%%%%%%%%
%% Final del capítulo de análisis del problema
%%%%%%%%%%%%%%%%%%%%%%%%%%%%%%%%%%%%%%%%%%%%%%%%%%%%%%%%%%%
